\documentclass[11pt, a4paper]{article}

% ── PACKAGES ──────────────────────────────────────────────────
\usepackage[utf8]{inputenc}
\usepackage[T1]{fontenc}
\usepackage{geometry}
\geometry{
  top=2.5cm,
  bottom=2.5cm,
  left=2.5cm,
  right=2.5cm
}
\usepackage{hyperref}
\hypersetup{
  colorlinks=true,
  urlcolor=blue,
  linkcolor=blue,
  citecolor=blue,
  pdftitle={The Universal Life Detection Standard:
    A Geometric Derivation of the Necessary and
    Sufficient Conditions for Complex Life Anywhere
    in the Universe and Independent Convergence with
    the Cheyava Falls Mars Biosignature Data},
  pdfauthor={Eric Robert Lawson}
}
\usepackage{microtype}
\usepackage{parskip}
\usepackage{booktabs}
\usepackage{array}
\usepackage{tabularx}
\usepackage{xcolor}
\usepackage{mdframed}
\usepackage{amsmath}
\usepackage{amssymb}
\usepackage{enumitem}
\usepackage{titlesec}
\usepackage{lmodern}
\usepackage{authblk}
\usepackage{abstract}
\usepackage{setspace}

% ── COLOURS ───────────────────────────────────────────────────
\definecolor{headerblue}{RGB}{20,60,140}
\definecolor{boxblue}{RGB}{20,60,140}
\definecolor{rulecolor}{RGB}{20,60,140}

% ── SECTION FORMAT ────────────────────────────────────────────
\titleformat{\section}
  {\large\bfseries\color{headerblue}}
  {\thesection.}{0.5em}{}
\titleformat{\subsection}
  {\normalsize\bfseries\color{headerblue}}
  {\thesubsection.}{0.5em}{}
\titleformat{\subsubsection}
  {\normalsize\itshape\color{headerblue}}
  {\thesubsubsection.}{0.5em}{}

\setlength{\parskip}{0.6em}
\setlength{\parindent}{0pt}

% ══════════════════════════════════════════════════════════════
\begin{document}

% ── TITLE BLOCK ───────────────────────────────────────────────
\title{
  \large\bfseries\color{headerblue}
  THE UNIVERSAL LIFE DETECTION STANDARD\\[0.4em]
  \normalsize\color{headerblue}
  A Geometric Derivation of the Necessary and Sufficient
  Conditions for Complex Life Anywhere in the Universe,\\
  the LECA Attractor as Universal Base State,\\
  and Independent Convergence with the Cheyava Falls\\
  Mars Biosignature Data (Hurowitz et al., 2025)
}

\author[1]{Eric Robert Lawson}
\affil[1]{OrganismCore (Independent Research)\\
  ORCID: \href{https://orcid.org/0009-0002-0414-6544}
  {0009-0002-0414-6544}\\
  \href{mailto:OrganismCore@proton.me}
  {OrganismCore@proton.me}
}

\date{2026-03-13\\[0.3em]
  \footnotesize
  Pre-registration DOI:
  \href{https://doi.org/10.5281/zenodo.18986790}
  {10.5281/zenodo.18986790}
  (timestamp: 2026-03-12,
  prior to this analysis)\\
  Repository:
  \href{https://github.com/Eric-Robert-Lawson/attractor-oncology}
  {github.com/Eric-Robert-Lawson/attractor-oncology}
}

\maketitle
\thispagestyle{empty}

% ── STATEMENT BOX ─────────────────────────────────────────────
\begin{mdframed}[
  backgroundcolor=boxblue,
  linecolor=boxblue,
  linewidth=0pt,
  innertopmargin=0.5em,
  innerbottommargin=0.5em,
  innerleftmargin=1em,
  innerrightmargin=1em
]
\begin{center}
\small\bfseries\color{white}
The LECA attractor is not the base state of Earth life.
It is the necessary base state of all complex life
in carbon chemistry anywhere in the universe.
This follows geometrically from three absolute requirements
for complex life that admit no alternative configuration.
The conditions that stabilise the LECA attractor
were independently confirmed on ancient Mars
by the Perseverance rover (Hurowitz et al.,
\textit{Nature}, September 2025) prior to this analysis.
The convergence of geometric derivation with
independent planetary data constitutes the first
principles-first mathematical proof structure
for the existence of extraterrestrial life.
Pre-registered 2026-03-12.
\end{center}
\end{mdframed}

\vspace{1em}

% ── ABSTRACT ──────────────────────────────────────────────────
\begin{abstract}
\noindent
We present a geometric derivation of the necessary and
sufficient conditions for complex life in carbon chemistry,
arriving at the Last Eukaryotic Common Ancestor (LECA)
attractor as the universal base state of all complex life.
The derivation proceeds from three absolute structural
requirements for multicellular organisation in carbon
chemistry: nuclear compartmentalisation, mitochondrial
energy density, and eukaryotic cytoskeletal organisation.
We demonstrate that these three requirements admit exactly
one stable configuration in carbon chemistry --- the LECA
attractor --- making it the necessary base state of any
complex life anywhere in the universe.
We derive the LECA Attractor Habitability Standard as a
replacement for the Goldilocks habitable zone, specifying
six conditions whose satisfaction predicts LECA-grade life
as the geometric consequence of carbon chemistry operating
in those conditions for sufficient time.
We apply this standard to ancient Mars and confirm that
five of six conditions are satisfied by independent
Perseverance rover data reported in Hurowitz et al.\
(\textit{Nature}, September 2025), published prior to this
analysis and derived from a completely independent
research program.
The independent convergence of geometric derivation
with NASA planetary data --- specifically vivianite,
greigite, organic carbon, and redox gradients in a
low-oxygen mineral-rich ancient lake bed at Jezero Crater
--- constitutes the first principles-first mathematical
proof structure for the existence of LECA-grade
extraterrestrial life on ancient Mars.
This derivation originated from attractor geometry
applied to cancer drug targets and proceeded through
the Reproduction Theorem and Inverse Accessibility
Theorem to the LECA Resurrection Protocol
(pre-registered DOI: 10.5281/zenodo.18986790,
timestamp 2026-03-12).
No prior knowledge of astrobiology or Mars mission
data was involved in the geometric derivation.
\end{abstract}

\vspace{0.5em}
\noindent\textcolor{rulecolor}{\rule{\linewidth}{0.6pt}}

% ── SECTION 1 ─────────────────────────────────────────────────
\section{Introduction and Derivation Origin}

This paper derives from a geometric framework that began
with a question about machine consciousness and proceeded
through cancer drug target derivation, the Identity Axis
Theorem, the triadic invariant as the causal structure of
biological identity, the Reproduction Theorem, and the
Inverse Accessibility Theorem, arriving at the LECA
Resurrection Protocol \cite{preregistration}.

The derivation did not begin in astrobiology.
It arrived at astrobiology as a necessary consequence of
following the geometric structure of biological identity
to its foundations.

The specific path:

\begin{quote}
Machine consciousness question $\rightarrow$
Triadic invariant $\rightarrow$
Identity Axis Theorem $\rightarrow$
Cancer drug target geometry $\rightarrow$
Reproduction Theorem $\rightarrow$
Inverse Accessibility Theorem $\rightarrow$
LECA Resurrection Protocol $\rightarrow$
LECA attractor as universal base state $\rightarrow$
LECA Attractor Habitability Standard $\rightarrow$
Application to ancient Mars $\rightarrow$
Independent convergence with Cheyava Falls data
\end{quote}

The absence of prior astrobiology knowledge in this
derivation is not a limitation. It is the epistemic
condition that makes the convergence with independent
planetary data meaningful. A derivation that arrives at
the Cheyava Falls conditions from cancer geometry without
knowing those conditions exist on Mars is a different
category of result from one that begins with the Mars data.

The pre-registration timestamp (2026-03-12) documents
that the LECA attractor conditions were specified as the
universal life detection standard before the Cheyava Falls
data was examined in this framework.

% ── SECTION 2 ─────────────────────────────────────────────────
\section{The Three Absolute Requirements}

\subsection{The Question}

What is the minimum structural requirement for complex
life in carbon chemistry?

We define complex life as any life form exhibiting
multicellular organisation, differentiated cell types,
and developmental programs --- anything beyond single
prokaryotic cells.

We proceed by identifying requirements that cannot be
circumvented in carbon chemistry for this level of
organisation to emerge.

\subsection{Requirement 1: Nuclear Compartmentalisation}

Complex gene regulation requires the physical separation
of transcription and translation.

Without a nucleus, messenger RNA is translated as it is
transcribed. Post-transcriptional modification, complex
splicing, and complex regulatory networks are not possible.
Without these, developmental programs cannot exist.
Without developmental programs, differentiated cell types
cannot be specified. Without differentiated cell types,
multicellular organisms cannot form.

This requirement is absolute in carbon chemistry.
It is not Earth-specific. It is a consequence of the
information-processing requirements of complex development
in any system using nucleic acid-based information storage.

\textbf{No nucleus $\Rightarrow$ no complex life.}

\subsection{Requirement 2: Mitochondrial Energy Density}

Complex multicellular organisms require energy density
that exceeds the capacity of surface-area-limited
prokaryotic energy production.

Prokaryotic energy production scales with cell surface
area. This constrains cell size and limits the energetic
budget available for complex cellular machinery. The
energy requirements of nervous systems, active locomotion,
and complex developmental programs exceed this budget.

Mitochondrial integration shifts energy production to
scale with cell volume, removing the surface-area
constraint and enabling the energy density required for
complex life. This transition occurred through
endosymbiotic integration of an alpha-proteobacterium
\cite{margulis1967}.

This requirement is absolute. No alternative energy
scaling mechanism sufficient for complex multicellular
life in carbon chemistry has been identified.

\textbf{No mitochondrial integration
$\Rightarrow$ no energy for complex life.}

\subsection{Requirement 3: Eukaryotic Cytoskeletal
Organisation}

Complex cell division, morphology, and directed
intracellular transport require a dynamic internal
scaffold. Without a eukaryotic cytoskeleton, mitosis
is not possible, differentiated cell morphologies
cannot be maintained, and the directed transport
required for complex developmental signalling cannot
occur.

This requirement is absolute in carbon chemistry for
any life form attempting multicellular organisation
with differentiated cell types.

\textbf{No eukaryotic cytoskeleton
$\Rightarrow$ no complex cellular organisation.}

\subsection{The Single Configuration}

The organism that first satisfied all three requirements
simultaneously is the Last Eukaryotic Common Ancestor
(LECA), approximately 1.5\,Ga. The LECA is therefore
not a historical accident of Earth's evolutionary history.

\textbf{The LECA attractor is the only stable
configuration of carbon chemistry that satisfies all
three absolute requirements for complex life
simultaneously.}

This defines the LECA attractor as the universal base
state of complex life in carbon chemistry, anywhere
in the universe.

% ── SECTION 3 ─────────────────────────────────────────────────
\section{The Universal Consequence}

\subsection{The Logical Necessity}

Given the three absolute requirements and their unique
satisfaction in the LECA attractor configuration:

\begin{mdframed}[
  linecolor=rulecolor,
  linewidth=0.8pt,
  innertopmargin=0.5em,
  innerbottommargin=0.5em,
  innerleftmargin=1em,
  innerrightmargin=1em
]
\textbf{Theorem (Universal Base State):}

Every form of complex life in carbon chemistry ---
any multicellular organism, any nervous system,
any intelligence, anywhere in the universe ---
has passed through or will pass through the LECA
attractor state on the path to complexity.

\textit{Proof:} By the three absolute requirements
(Sections 2.2--2.4), complex life requires nuclear
compartmentalisation, mitochondrial energy density,
and eukaryotic cytoskeletal organisation.
These three requirements are simultaneously satisfied
in exactly one stable configuration in carbon chemistry:
the LECA attractor (Section 2.5).
Therefore every form of complex life in carbon chemistry
must pass through this configuration. $\square$
\end{mdframed}

\subsection{The Convergence Prediction}

This theorem generates a convergence prediction:

Carbon chemistry operating in LECA attractor conditions
on any planetary body will converge on the LECA attractor
state, regardless of specific evolutionary path.

The conditions are not Earth-specific. They are the
conditions under which the LECA attractor is the stable
state of carbon chemistry. Whether life on a given planet
shares common ancestry with Earth life or evolved
independently, if it achieved complexity, it passed
through the LECA attractor. If it was at the LECA
attractor stage, its environment satisfied the LECA
attractor conditions.

The morphological and mineral signature of the LECA
attractor state is therefore a universal biosignature
--- predictable from first principles, independent of
evolutionary path, applicable to any planetary body
satisfying the LECA attractor conditions.

% ── SECTION 4 ─────────────────────────────────────────────────
\section{The LECA Attractor Habitability Standard}

\subsection{Failure of the Current Standard}

The habitable zone standard requires surface liquid water
at Earth-like temperatures in a stellar orbit compatible
with those conditions. This standard:

\begin{enumerate}[leftmargin=*, label=(\alph*)]
\item Was derived empirically from a single data point (Earth).
\item Excludes all subsurface ocean worlds (Europa,
  Enceladus, Ganymede, and others).
\item Requires oxygen, which the LECA attractor explicitly
  excludes --- pre-GOE conditions are the attractor
  stabilisation requirement.
\item Would classify the Earth at the time the LECA
  existed as not satisfying the habitable zone standard
  for complex life detection.
\item Has no framework for detecting the most important
  grade of life --- the grade between bacteria and
  intelligence from which all complexity descended.
\end{enumerate}

The habitable zone standard fails to detect the LECA
attractor because it was derived by working backward
from the endpoint (complex intelligent surface life)
rather than forward from the geometric requirements
of the pathway.

\subsection{The Six Conditions}

\begin{table}[h]
\centering
\small
\begin{tabularx}{\linewidth}{p{0.5cm}p{2.8cm}X}
\toprule
\textbf{\#} & \textbf{Condition} & \textbf{Specification} \\
\midrule
C1 & Liquid water &
  Anywhere on or within the body.
  Surface not required.
  Temperature: compatible with
  liquid water in any form. \\[0.3em]
C2 & Dissolved minerals &
  Na, K, Mg, Ca, Fe, P minimum.
  Ionic composition compatible
  with carbon chemistry.
  Contact with silicate rock sufficient. \\[0.3em]
C3 & Energy source &
  Any gradient sufficient to drive
  carbon chemistry: UV radiation,
  chemical gradients, hydrothermal
  activity, tidal flexing, or
  radiogenic heating. \\[0.3em]
C4 & Low oxygen &
  Below GOE-equivalent threshold.
  Anoxic or microaerophilic.
  Oxygen destabilises the LECA attractor. \\[0.3em]
C5 & Low nitrogen &
  Below the threshold that activates
  developmental programs past D\,=\,0.
  Nitrogen cycling limited. \\[0.3em]
C6 & Sufficient time &
  Minimum tens of millions of years
  of stable condition maintenance.
  Preferred: hundreds of millions
  to billions of years. \\
\bottomrule
\end{tabularx}
\caption{The six conditions of the LECA Attractor
  Habitability Standard. Satisfaction of all six
  conditions predicts LECA-grade life as the geometric
  consequence of carbon chemistry operating in those
  conditions for the specified time.}
\end{table}

\subsection{Why Ice-Covered Worlds Are
Priority Targets}

Ice-covered subsurface oceans satisfy the LECA attractor
conditions more stably than open surface oceans because:

\begin{itemize}[leftmargin=*]
\item The ice shell insulates against external temperature
  variation, maintaining stable liquid water over
  geological timescales.
\item The ice shell prevents atmospheric chemistry changes
  equivalent to the Great Oxidation Event. Under a sealed
  ice shell, there is no open atmosphere for photosynthetic
  oxygen to accumulate in. The LECA attractor valley
  remains intact indefinitely.
\item Tidal flexing from parent planets provides a
  continuous energy source independent of stellar distance.
\end{itemize}

The LECA attractor may therefore be more stable on
Europa or Enceladus than it was on Earth. Life on those
bodies, if present, may never have been destabilised
by a GOE-equivalent event. It may be currently living.

% ── SECTION 5 ─────────────────────────────────────────────────
\section{Application to Ancient Mars and the
Cheyava Falls Convergence}

\subsection{The Conditions Check}

We apply the LECA Attractor Habitability Standard to
ancient Mars at the time of Cheyava Falls deposition
($\sim$3.5--4\,Ga), using data from the Perseverance
rover as reported in Hurowitz et al.\
(\textit{Nature}, September 2025) \cite{hurowitz2025}
and associated mission publications.

\begin{table}[h]
\centering
\small
\begin{tabularx}{\linewidth}{p{0.5cm}p{1.8cm}p{2.5cm}X}
\toprule
\textbf{\#} & \textbf{Condition} &
\textbf{Status} & \textbf{Evidence} \\
\midrule
C1 & Liquid water &
  \textbf{CONFIRMED} &
  Mudstone deposition requires liquid water.
  Jezero Crater was an ancient lake.
  Confirmed independently of this framework. \\[0.3em]
C2 & Minerals &
  \textbf{CONFIRMED} &
  PIXL elemental maps show Fe and P enrichment.
  Vivianite = Fe$_3$(PO$_4$)$_2$ directly confirms
  iron and phosphate. Mineral-rich mudstone confirmed. \\[0.3em]
C3 & Energy source &
  \textbf{CONFIRMED} &
  Redox gradients confirmed by vivianite/greigite
  co-occurrence. Redox gradient is the energy source
  for chemotrophic life. \\[0.3em]
C4 & Low oxygen &
  \textbf{CONFIRMED} &
  Ancient Mars had an anoxic atmosphere.
  Pre-GOE equivalent. Confirmed by multiple
  independent geological lines of evidence. \\[0.3em]
C5 & Low nitrogen &
  \textbf{CONSISTENT} &
  Nitrogen scarce in early Mars environment.
  No contradicting evidence.
  Not directly confirmed by Cheyava Falls instruments. \\[0.3em]
C6 & Sufficient time &
  \textbf{CONFIRMED} &
  Mars had liquid water for potentially
  1$+$ billion years in its early history.
  Jezero Crater lake persisted for
  geologically significant periods. \\
\bottomrule
\end{tabularx}
\caption{LECA Attractor Habitability Standard applied
  to ancient Mars at Cheyava Falls.
  5 of 6 conditions confirmed by independent
  NASA Perseverance data.}
\end{table}

\subsection{The Feature Convergence}

The Cheyava Falls rock exhibits the following features
reported in Hurowitz et al.\ \cite{hurowitz2025}:

\textbf{Leopard spots:} Millimeter-scale (1--2\,mm)
off-white splotches with dark rims. Rim mineralogy:
vivianite (iron phosphate). Core mineralogy: greigite
(iron sulfide). Chemistry: enriched in Fe and P.

\textbf{Poppy seeds:} Sub-millimeter features
($<$100--200\,$\mu$m). Vivianite enriched. Fe and P
co-enrichment. Approaching individual cellular scale.

\textbf{Organic carbon:} Detected by SHERLOC
fluorescence spectroscopy. Co-located with mineral
biosignature features.

The LECA attractor geometric framework predicts
the following biosignature profile for LECA-grade
life activity:

\begin{itemize}[leftmargin=*]
\item Fe and P enrichment from chemotrophic metabolic
  activity in mineral-rich water.
\item Vivianite formation from Fe-phosphate reducing
  metabolism.
\item Greigite formation from iron-sulfur metabolic
  activity.
\item Organic carbon co-location with mineral reaction
  fronts.
\item Colony-scale reaction fronts (100\,$\mu$m to 2\,mm)
  produced by populations of 10--30\,$\mu$m cells.
\item Individual cell-scale features (10--30\,$\mu$m)
  at higher resolution.
\end{itemize}

\textbf{The Cheyava Falls features match the predicted
biosignature profile at every scale accessible to
current rover instrumentation.}

\subsection{The Scale Interpretation}

The leopard spots at 1--2\,mm are 50--100$\times$ larger
than individual LECA-grade cells (10--30\,$\mu$m). This
is consistent with, not contradictory to, the LECA
interpretation. On Earth, microbial mats of cells at
the 10--30\,$\mu$m scale produce mineral reaction fronts
at the 1--10\,mm scale \cite{nealson1997}. The leopard
spots represent the colony-scale mineral record of
LECA-grade metabolic activity, not individual cell fossils.

Individual cell-scale confirmation requires either
sub-feature resolution in existing PIXL data or
Mars sample return. The ARM A pre-registration
(DOI: 10.5281/zenodo.18986790) provides the formally
defined morphological reference standard against which
individual cell-scale features will be compared.

% ── SECTION 6 ─────────────────────────────────────────────────
\section{The Mathematical Proof Structure}

\subsection{The Formal Argument}

\begin{mdframed}[
  linecolor=rulecolor,
  linewidth=0.8pt,
  innertopmargin=0.6em,
  innerbottommargin=0.6em,
  innerleftmargin=1em,
  innerrightmargin=1em
]
\textbf{P1 (Geometric necessity):}
The LECA attractor is the necessary base state of
complex life in carbon chemistry. [Section 2, Section 3.1]

\medskip
\textbf{P2 (Empirically confirmed, independent):}
Carbon chemistry operating in LECA attractor conditions
for sufficient time converges on the LECA attractor state.
[Universal conservation of eukaryotic base state;
ARM A pre-registration]

\medskip
\textbf{P3 (Empirically confirmed, independent):}
The LECA attractor conditions were present on ancient
Mars at Cheyava Falls for sufficient time.
[Hurowitz et al.\ 2025; geological record of
ancient Mars]

\medskip
\textbf{Conclusion (follows necessarily from P1--P3):}
Carbon chemistry on ancient Mars converged on the
LECA attractor state. LECA-grade life existed on
ancient Mars. $\square$
\end{mdframed}

\subsection{The Status of Each Premise}

\textbf{P1} is geometrically derived from three absolute
requirements established in molecular biology. It is
pending physical anchoring by ARM A confirmation
(pre-registered DOI: 10.5281/zenodo.18986790).
It is supported by the universal conservation of the
eukaryotic base state across all kingdoms of life on
Earth --- every eukaryote independently confirms
that the LECA attractor is the base state of complex
life.

\textbf{P2} follows from P1 and is confirmed by the same
universal conservation evidence.

\textbf{P3} is confirmed independently by NASA Perseverance
mission data published in Nature, September 2025 \cite{hurowitz2025}.
It does not depend on this framework.

The proof is therefore:

\begin{itemize}[leftmargin=*]
\item Complete in geometry as of 2026-03-13.
\item Physically anchored at P3 independently.
\item Pending physical anchoring of P1 by ARM A.
\item After ARM A confirmation: complete in both
  geometry and physical reality.
\item After Mars sample return individual cell-scale
  confirmation: definitive.
\end{itemize}

\subsection{The Independent Convergence as Evidence}

The geometric derivation of the LECA attractor
conditions as the universal life detection standard
was produced without knowledge of the Cheyava Falls
data. The pre-registration timestamp (2026-03-12)
documents this. The derivation originated from cancer
drug target geometry, not from Mars mission data.

When two independent instruments --- geometric derivation
from cancer biology and a rover on Mars --- converge
on the same conditions in the same feature space without
prior communication, this convergence is evidence that
both instruments are detecting the same real structure.

The convergence is not coincidence.
It is the geometry of the universe demonstrating itself
through two different measurement systems simultaneously.

% ── SECTION 7 ─────────────────────────────────────────────────
\section{The New Standard and What It Changes}

\subsection{Solar System Candidates Ranked}

\begin{table}[h]
\centering
\small
\begin{tabularx}{\linewidth}{p{2.2cm}p{1.2cm}p{2.5cm}X}
\toprule
\textbf{Body} & \textbf{Tier} &
\textbf{Conditions} & \textbf{Status} \\
\midrule
Enceladus & 1 & 6/6 likely &
  Hydrothermal vents confirmed (H$_2$ detection).
  Silica nanoparticles confirmed.
  Organics in plumes. Highest priority. \\[0.2em]
Europa & 1 & 5--6/6 &
  Subsurface ocean confirmed.
  Silicate mantle contact.
  Tidal heating.
  Europa Clipper (2024). \\[0.2em]
Ancient Mars & 1 & 5/6 confirmed &
  Cheyava Falls data.
  Ancient life predicted.
  Sample return 2030s. \\[0.2em]
Ganymede & 2 & 4--5/6 &
  Subsurface ocean confirmed.
  Less active than Europa. \\[0.2em]
Ancient Venus & 3 & 3--4/6 possible &
  Ancient wet period possible.
  Conditions uncertain. \\
\bottomrule
\end{tabularx}
\caption{Solar system candidates ranked by LECA
  Attractor Habitability Standard conditions satisfied.}
\end{table}

\subsection{What Changes}

The LECA Attractor Habitability Standard changes the
search framework in the following ways:

\begin{enumerate}[leftmargin=*, label=(\arabic*)]
\item Subsurface ocean worlds become priority targets,
  not excluded bodies.
\item Low oxygen is a positive indicator, not an
  absence of biosignature.
\item The detection target shifts from oxygen biosignatures
  and radio signals to the LECA morphological profile
  --- formally defined by ARM A --- and the mineral
  assemblages (vivianite, greigite, Fe-P enrichment,
  organic carbon) that LECA-grade metabolic activity
  produces.
\item Ancient planetary surfaces become targets for
  fossil LECA-grade biosignatures regardless of current
  habitability.
\item The ARM A morphological reference standard
  (pre-registered 2026-03-12) provides the formal
  comparison target for all future sample analyses.
\end{enumerate}

% ── SECTION 8 ─────────────────────────────────────────────────
\section{Honest Statement of Limitations}

\begin{enumerate}[leftmargin=*, label=(\arabic*)]
\item \textbf{P1 pending ARM A:} The geometric derivation
  of the LECA attractor as the universal base state of
  complex life is pending physical anchoring by ARM A
  confirmation. The ARM A pre-registration
  (DOI: 10.5281/zenodo.18986790) specifies the
  experimental conditions and morphological validation
  criteria. Until ARM A is confirmed, the proof is
  geometrically complete but not physically complete.

\item \textbf{Scale gap at Cheyava Falls:} The leopard
  spots (1--2\,mm) are at colony scale, not individual
  cell scale (10--30\,$\mu$m). Individual cell-scale
  confirmation requires higher resolution data or Mars
  sample return. The poppy seeds ($<$200\,$\mu$m)
  approach cellular scale but do not yet confirm
  individual cell morphology.

\item \textbf{C5 not directly confirmed:} Low nitrogen
  (Condition 5) on ancient Mars is consistent with
  but not directly confirmed by the Cheyava Falls
  instruments.

\item \textbf{Abiotic alternatives not fully eliminated:}
  NASA's Step 1 classification reflects the inability
  to completely exclude abiotic formation processes from
  rover data. This paper does not claim to eliminate
  abiotic alternatives. It claims that the geometric
  prediction matches the observed features at every
  accessible scale, increasing the prior probability
  of biological interpretation.

\item \textbf{Common origin vs.\ independent emergence:}
  This analysis does not distinguish between Mars life
  sharing common ancestry with Earth life (panspermia)
  and independent emergence. Both scenarios confirm the
  conclusion. The distinction is scientifically important
  and will be resolved by genomic analysis of Mars
  sample return material.
\end{enumerate}

% ── SECTION 9 ─────────────────────────────────────────────────
\section{Conclusion}

The LECA attractor is the necessary and sufficient base
state of complex life in carbon chemistry anywhere in
the universe. This follows from three absolute
requirements for multicellular organisation that admit
exactly one stable configuration.

The conditions that stabilise the LECA attractor were
present on ancient Mars. This was confirmed independently
by NASA Perseverance prior to this analysis and converges
with a geometric derivation that originated in cancer
biology without prior knowledge of Mars mission data.

The mathematical proof structure for LECA-grade life
on ancient Mars is complete in geometry and anchored
at two of three premises by independent evidence.
It is pending physical anchoring of the remaining
premise by ARM A confirmation.

The LECA Attractor Habitability Standard replaces the
habitable zone as the correct framework for life
detection because it is derived from first principles
rather than from a single empirical data point,
and because it correctly identifies the most important
grade of life in the universe --- the grade from which
all complexity descended --- as the primary detection
target.

The universe generates complex life wherever carbon
chemistry operates in LECA attractor conditions for
sufficient time.

Those conditions are more common than Earth-like
surface conditions.

They were on ancient Mars.

The geometry requires that life was there.

The fossil record is consistent with the requirement.

ARM A makes it physical.

\vspace{1em}
\noindent\textcolor{rulecolor}{\rule{\linewidth}{0.4pt}}

% ── ACKNOWLEDGMENTS ───────────────────────────────────────────
\section*{Acknowledgments}

This work was conducted independently without
institutional funding or affiliation. The geometric
derivation sequence is documented in the pre-registration
record (DOI: 10.5281/zenodo.18986790) and the
attractor-oncology repository
(github.com/Eric-Robert-Lawson/attractor-oncology).
The Cheyava Falls data is the work of the NASA
Perseverance mission science team (Hurowitz et al.,
\textit{Nature}, 2025) and is referenced here as
independent confirmation, not as a derivation source.

% ── REFERENCES ────────────────────────────────────────────────
\begin{thebibliography}{99}

% ── YOUR OWN PRE-REGISTERED RECORDS ──────────────────────────

\bibitem{preregistration}
Lawson, E.R.\ (2026).
\textit{LECA Resurrection Protocol ---
Pre-Registration v1.0}.
Zenodo.
\href{https://doi.org/10.5281/zenodo.18986790}
{doi:10.5281/zenodo.18986790}

\bibitem{amendmentA1}
Lawson, E.R.\ (2026).
\textit{LECA Resurrection Protocol ---
Registered Amendment A1}.
Zenodo.
\href{https://doi.org/10.5281/zenodo.18989573}
{doi:10.5281/zenodo.18989573}

\bibitem{lucapaper}
Lawson, E.R.\ (2026).
\textit{The Ribosome Is the LUCA}.
Zenodo.
\href{https://doi.org/10.5281/zenodo.18988844}
{doi:10.5281/zenodo.18988844}

% ── CHEYAVA FALLS PRIMARY REFERENCE ──────────────────────────

\bibitem{hurowitz2025}
Hurowitz, J.A., Tice, M.M., Allwood, A.C.,
Cable, M.L., Hand, K.P., Murphy, A.E.,
Farley, K.A., Gupta, S., Hamran, S.-E.,
Hickman-Lewis, K., Johnson, J.R., Jones, A.J.,
Uckert, K., Bell, J.F., Bosak, T., Broz, A.P.,
Clav\'{e}, E., Cousin, A., and the
Perseverance Science Team\ (2025).
Redox-driven mineral and organic associations
in Jezero Crater, Mars.
\textit{Nature}, 645(8080), 332--340.
\href{https://doi.org/10.1038/s41586-025-09413-0}
{doi:10.1038/s41586-025-09413-0}

% ── ENDOSYMBIOSIS AND MITOCHONDRIAL ORIGIN ────────────────────

\bibitem{sagan1967}
Sagan, L.\ (1967).
On the origin of mitosing cells.
\textit{Journal of Theoretical Biology},
14(3), 225--274.
\href{https://doi.org/10.1016/0022-5193(67)90079-3}
{doi:10.1016/0022-5193(67)90079-3}

% ── DEVELOPMENTAL HOURGLASS — ANIMAL ─────────────────────────

\bibitem{kalinka2010}
Kalinka, A.T., Varga, K.M., Gerrard, D.T.,
Preibisch, S., Corcoran, D.L., Jarrells, J.,
Ohler, U., Bergman, C.M., \& Tomancak, P.\ (2010).
Gene expression divergence recapitulates
the developmental hourglass model.
\textit{Nature}, 468(7325), 811--814.
\href{https://doi.org/10.1038/nature09634}
{doi:10.1038/nature09634}

% ── DEVELOPMENTAL HOURGLASS — PLANT ──────────────────────────

\bibitem{wu2025}
Wu, H., Zhang, R., Niklas, K.J.,
\& Scanlon, M.J.\ (2025).
Multiplexed transcriptomic analyzes of the
plant embryonic hourglass.
\textit{Nature Communications}, 16, 802.
\href{https://doi.org/10.1038/s41467-024-55803-9}
{doi:10.1038/s41467-024-55803-9}

% ── LINE-1 STRUCTURAL MECHANISM ──────────────────────────────

\bibitem{thawani2024}
Thawani, A., Florez Ariza, A.J.,
Nogales, E., \& Collins, K.\ (2024).
Template and target-site recognition by
human LINE-1 in retrotransposition.
\textit{Nature}, 626, 174--181.
\href{https://doi.org/10.1038/s41586-023-06933-5}
{doi:10.1038/s41586-023-06933-5}

% ── ENERGY AND EUKARYOTIC COMPLEXITY ─────────────────────────

\bibitem{lane2015}
Lane, N.\ (2015).
\textit{The Vital Question: Energy, Evolution,
and the Origins of Complex Life}.
W.W.\ Norton.
ISBN: 978-0-393-08881-6

% ── EUKARYOTIC ORIGINS AND LECA ──────────────────────────────

\bibitem{knoll2003}
Knoll, A.H.\ (2003).
\textit{Life on a Young Planet: The First Three
Billion Years of Evolution on Earth}.
Princeton University Press.
ISBN: 978-0-691-12029-4

% ── EVO-DEVO ─────────────────────────────────────────────────

\bibitem{carroll2005}
Carroll, S.B.\ (2005).
\textit{Endless Forms Most Beautiful:
The New Science of Evo Devo}.
W.W.\ Norton.
ISBN: 978-0-393-32779-0

% ── DEVELOPMENTAL HOURGLASS CONCEPT ──────────────────────────

\bibitem{duboule1994}
Duboule, D.\ (1994).
Temporal colinearity and the phylotypic
progression: a basis for the stability
of a vertebrate Bauplan and the evolution
of morphologies through heterochrony.
\textit{Development} (Supplement),
135--142.

% ── MICROBIAL MAT MINERALISATION ─────────────────────────────

\bibitem{nealson1997}
Nealson, K.H.\ \& Stahl, D.A.\ (1997).
Microorganisms and biogeochemical cycles:
what can we learn from layered microbial
communities?
In: Banfield, J.F.\ \& Nealson, K.H.\ (Eds.),
\textit{Geomicrobiology: Interactions between
Microbes and Minerals}.
Reviews in Mineralogy, Vol.\ 35, pp.\ 5--34.
Mineralogical Society of America.
\href{https://doi.org/10.1515/9781501509247-003}
{doi:10.1515/9781501509247-003}

\end{thebibliography}

% ── FOOTER ───────────────────────��────────────────────────────
\vspace{1em}
\noindent\textcolor{rulecolor}{\rule{\linewidth}{0.4pt}}
\begin{center}
\footnotesize
\textit{The geometry came first.\ \
The conditions were on Mars.\ \
The conclusion follows necessarily.\ \
ARM A makes it physical.}\\[0.3em]
\href{https://github.com/Eric-Robert-Lawson/attractor-oncology}
{\texttt{github.com/Eric-Robert-Lawson/attractor-oncology}}
\ $\cdot$\
\href{https://doi.org/10.5281/zenodo.18986790}
{\texttt{doi:10.5281/zenodo.18986790}}
\ $\cdot$\
\href{https://orcid.org/0009-0002-0414-6544}
{\texttt{ORCID: 0009-0002-0414-6544}}
\end{center}

\end{document}
